% ECONOMICS_PROBLEM: {"id":"D82-92","title":"Multidimensional screening","jel_code":"D82","source_id":"OP-0006"}
% FIRST_POSTED: 2026-10-10
% ENTRY_KIND: research_attempt
% ATTEMPT_STATUS: solved
\documentclass[11pt,a4paper]{article}
\usepackage[margin=25mm]{geometry}
\usepackage{amsmath,amssymb,amsthm}
\usepackage[hidelinks]{hyperref}
\hypersetup{pdftitle={Multidimensional Screening: Global Certificates and Exact Geometry},pdfauthor={Ji Ho Bae}}
\setlength{\parindent}{0pt}
\setlength{\parskip}{4pt}
\setlength{\emergencystretch}{2em}
\allowdisplaybreaks[2]
\numberwithin{equation}{section}
\newtheorem{theorem}{Theorem}
\newtheorem{proposition}[theorem]{Proposition}
\newtheorem{lemma}[theorem]{Lemma}
\newtheorem{corollary}[theorem]{Corollary}
\newcommand{\R}{\mathbb R}
\newcommand{\E}{\mathbb E}
\newcommand{\Prob}{\mathbb P}
\newcommand{\norm}[1]{\left\lVert #1\right\rVert}
\newcommand{\T}{\mathsf T}
\newcommand{\ind}{\mathbf1}
\DeclareMathOperator{\diam}{diam}
\DeclareMathOperator{\rank}{rank}
\DeclareMathOperator{\vol}{vol}
\DeclareMathOperator{\supp}{supp}
\title{Multidimensional Screening:\\Global Certificates and Exact Geometry}
\author{Ji Ho Bae\thanks{Prepared with GPT Codex assistance; the exact serving revision is unavailable. All proofs are hand arguments. No Lean formalization, external peer review, or novelty claim is asserted.}}
\date{10 October 2026}
\begin{document}
\maketitle
\begin{center}
\small\textbf{D82-92 \quad JEL D82 \quad Source OP-0006}\\
Complete hand proof with internal mathematical and scope review.
\end{center}
\begin{abstract}
For compact convex type and quantity sets, a continuously positive type
density, and a continuous convex production cost, we prove existence of
an optimal globally incentive-compatible and individually rational
mechanism, retaining feasibility at every boundary type. Strict cost
convexity yields the appropriate almost-everywhere uniqueness, while an
explicit example shows nonuniqueness otherwise. A general finite-menu
construction covers signed costs, non-Lipschitz costs, and quantity sets
excluding zero. Under explicit effective representations it produces a
certified global objective gap. A sufficient measure-flow dual witness
provides another global verifier. Finally, a full-dimensional radial
class admits an exact optimizer, an all-competitor deficit identity,
precise exclusion and served rank-deficient bunching conditions, and
certified rational hinge computation. These structural geometry results
are not a classification for every arbitrary density and cost.
\end{abstract}

\section{Primitives, scope, and global implementation}\label{sec:primitives}
Let \(d\ge2\). Let \(\Theta\subset\R^d\) be compact and convex and carry
a continuously positive probability density \(f\) relative to
\(d\)-dimensional Lebesgue measure. Write
\(\mu(d\theta)=f(\theta)\,d\theta\).
Consequently \(\Theta\) has positive volume and nonempty interior.
Let \(Q\subset\R_+^d\) be nonempty, compact, and convex, and let
\(C:Q\to\R\) be continuous and convex.
If \(Q\) is empty, the problem is infeasible.
We impose no zero-allocation assumption, nonnegativity or Lipschitz
assumption on cost, smoothness assumption, or rectangular type domain.

A mechanism consists of measurable \(q:\Theta\to Q\) and
\(t:\Theta\to\R\), with utility
\(U(\theta)=\theta\cdot q(\theta)-t(\theta)\).
It is globally IC and IR when
\[
 U(\theta)\ge\theta\cdot q(\phi)-t(\phi)
 \quad(\theta,\phi\in\Theta),\qquad U(\theta)\ge0.
\]
All these inequalities, including those at the boundary, are required.
Its profit is
\[
 J(U,q)=\int_\Theta[\theta\cdot q(\theta)-U(\theta)-C(q(\theta))]\,d\mu,
 \qquad V^*=\sup_{\mathrm{IC,IR}}J.
\]
Choose finite certified bounds
\[
 H\ge\max\{1,\sup_\Theta\norm\theta\},\quad
 R_Q\ge\max\{1,\sup_Q\norm q\},\quad
 D\ge\diam(\Theta),\quad A\ge\sup_Q|C|,
\]
and put \(K=(H+D)R_Q\), \(P=K+A\).
Let \(\omega\) be a nondecreasing upper cost modulus:
\[
 |C(q)-C(q')|\le\omega(\norm{q-q'}),\qquad
 \omega(r)\longrightarrow0\quad(r\downarrow0).
\]

\begin{lemma}[Global potential equivalence and normalization]\label{lem:normal}
Every IC mechanism has a convex \(R_Q\)-Lipschitz utility, and its
assigned allocation is a relative subgradient in \(Q\) at every type.
Conversely, these support inequalities and \(U\ge0\) imply global
IC/IR for \(t=\theta\cdot q-U\).
Every feasible mechanism can be normalized without reducing profit so that
\begin{equation}\label{eq:normalbounds}
 \min_\Theta U=0,\qquad
 0\le U\le R_QD,\qquad |t|\le K,\qquad |t-C(q)|\le P.
\end{equation}
\end{lemma}
\begin{proof}
Two IC inequalities give
\[
 (\theta-\phi)\cdot q(\phi)
 \le U(\theta)-U(\phi)
 \le(\theta-\phi)\cdot q(\theta).
\]
These imply the Lipschitz bound. Also
\[
 U(\theta)=\sup_{\phi\in\Theta}\{\theta\cdot q(\phi)-t(\phi)\},
\]
so \(U\) is convex and
\(U(\phi)\ge U(\theta)+q(\theta)\cdot(\phi-\theta)\).
The converse follows by substituting the payment formula.
Continuity gives \(m=\min_\Theta U\ge0\). Raising every payment by
\(m\) preserves IC/IR and raises expected profit by \(m\).
The remaining bounds follow from the Lipschitz estimate and
\(t=\theta\cdot q-U\).
\end{proof}

Feasibility is nonempty: fix \(q_0\in Q\) and charge the constant
\(t_0=\min_\Theta\theta\cdot q_0\).
Together with \eqref{eq:normalbounds}, this proves \(V^*\) is finite.
This contract is used only to prove nonemptiness; it will not be inserted
into an existing menu.
If \(0\notin Q\), no-trade exclusion is unavailable, and \(U=0\) must
not be renamed exclusion. If \(0\in Q\) and a normalized mechanism assigns
\(q(\theta)=0\), the support inequality makes that type a global utility
minimizer, so \(U(\theta)=t(\theta)=0\).

\section{Existence and the precise uniqueness statement}\label{sec:existence}
\begin{theorem}[Compact existence]\label{thm:existence}
The supremum \(V^*\) is attained under the primitives of
Section~\ref{sec:primitives}. The attaining mechanism is measurable
and satisfies global IC/IR at every type, including the boundary.
\end{theorem}
\begin{proof}
Take normalized feasible mechanisms \((U_n,q_n,t_n)\) whose profits
converge to \(V^*\). Their utilities are uniformly bounded and
\(R_Q\)-Lipschitz. To obtain a uniformly convergent subsequence, choose
finite \(1/j\)-nets in \(\Theta\), diagonalize on their countable union,
and use the common Lipschitz bound to turn convergence on a finite net
into the uniform Cauchy property. Relabel the subsequence.
Its limit \(U\) is convex, \(R_Q\)-Lipschitz, nonnegative, and has
minimum zero.

We construct a measurable pointwise cluster selection of \(q_n\).
Let \(a_1(\theta)=\liminf_n q_n^1(\theta)\). Recursively choose the first
index \(n_j(\theta)>n_{j-1}(\theta)\) with
\[
 |q_{n_j(\theta)}^1(\theta)-a_1(\theta)|<1/j.
\]
Such indices exist arbitrarily far out by the definition of the
bounded-coordinate lower limit. The first-index choice among countably
many measurable conditions is measurable. Take the lower limit of the
second coordinate along this measurable subsequence and extract again,
preserving convergence of the first coordinate. After \(d\) finite
extractions all coordinates converge along a measurable subsequence.
Its limit \(q(\theta)\) is measurable and belongs to compact \(Q\).
Endpoint lower limits cause no problem because every coordinate is bounded.

For each fixed \(\theta\), pass along its selected subsequence in
\[
 U_n(\phi)\ge U_n(\theta)+q_n(\theta)\cdot(\phi-\theta).
\]
Uniform convergence yields the same inequality for \(U,q\), for every
\(\phi\), with the same selected \(q(\theta)\).
Thus \(q(\theta)\) is a relative subgradient in \(Q\) at every type.
The payment \(t(\theta)=\theta\cdot q(\theta)-U(\theta)\) is measurable
and bounded, and the resulting mechanism is globally IC/IR.
No exceptional set was discarded in constructing this mechanism.

It remains to show convergence of profit. We record a direct
almost-everywhere differentiability argument. On each coordinate line,
the two one-sided derivatives of a convex function differ at only
countably many points: the nonempty intervals between them are disjoint,
and each contains a distinct rational. Fubini gives simultaneous
coordinate derivatives almost everywhere in the interior.
At such a point every subgradient has those derivatives as its
coordinates, and is therefore unique. Bounded subgradients at nearby
points have cluster points only in that singleton, by passage in the
support inequalities. Sandwiching utility increments between the two
supporting planes proves differentiability there.
Subgradients at every point were supplied above and are bounded by \(R_Q\).

The boundary has zero volume. Choose an interior ball \(B(z,r)\subset\Theta\).
The homothetic copy \(z+(1-\varepsilon)(\Theta-z)\) lies in the interior:
each point in it has an \(\varepsilon r\)-ball inside \(\Theta\).
Thus
\[
 \vol(\partial\Theta)
 \le[1-(1-\varepsilon)^d]\vol(\Theta)\longrightarrow0.
\]
At an interior differentiability point of \(U\), every cluster point of
the original \(q_n(\theta)\) is a subgradient of \(U\), hence equals
\(\nabla U(\theta)\). Compactness of \(Q\) forces convergence of the
entire sequence there to \(q(\theta)=\nabla U(\theta)\).
Continuity of \(C\), \eqref{eq:normalbounds}, and dominated convergence
now give convergence of
\[
 \int_\Theta[\theta\cdot q_n-U_n-C(q_n)]\,d\mu
\]
to the profit of the constructed mechanism. That profit is \(V^*\).
\end{proof}

\begin{theorem}[Strict-cost uniqueness]\label{thm:unique}
If \(C\) is strictly convex on \(Q\), every optimal mechanism has the
same utility at every type and the same allocation and payment almost
everywhere. Pointwise uniqueness of boundary allocations is not asserted.
\end{theorem}
\begin{proof}
If two optimal allocations differ on a set of positive probability,
their midpoint mechanism is globally IC/IR by averaging support
inequalities and using convexity of \(Q\). Strict convexity of cost
strictly increases its integrated profit, a contradiction.
The allocations therefore agree almost everywhere, as do the gradients
of the two Lipschitz utilities on the connected interior.
Their difference is constant: on each interior ball mollify the
difference; its gradient is the convolution of its zero weak gradient,
so it is constant on smaller balls. Pass to the uniform limit and join
overlapping balls. The weak-gradient statement follows also from
absolute continuity on coordinate lines. Continuity extends the
constant to the boundary. Optimal utilities have minimum zero by
Lemma~\ref{lem:normal}, so the constant is zero. Payments agree
wherever the allocations do.
\end{proof}

\subsection{A concrete nonunique optimum}\label{sec:nonunique}
Let
\[
 \Theta=[1,3]\times[0,1]^{d-1},\qquad
 Q=\{ze_1:0\le z\le1\},\qquad C=0.
\]
Use a product density, uniform in the last coordinates, with
\[
 f_1(x)=
 \begin{cases}x^{-2},&1\le x\le2,\\
 \frac14+\frac{x-2}{2},&2\le x\le3.\end{cases}
\]
Both pieces have integral \(1/2\), and they join continuously and
positively at \(2\). Its tail is
\[
 S(x)=
 \begin{cases}x^{-1},&1\le x\le2,\\
 \frac12-\frac{x-2}{4}-\frac{(x-2)^2}{4},&2\le x\le3.
 \end{cases}
\]
Applying support inequalities between types with the same \(x\) shows
that every feasible utility depends only on \(x\).
Its scalar allocation equals a nondecreasing derivative \(z(x)\)
almost everywhere. Normalize \(U(1)=0\). Fubini gives, for every mechanism,
\[
 J=\int_1^3[xf_1(x)-S(x)]z(x)\,dx.
\]
The coefficient is zero on \([1,2]\). It is zero at \(2\) and strictly
increasing on \((2,3)\), since its derivative is \(2f_1+xf_1'>0\).
As \(0\le z\le1\), the resulting upper bound equals \(1\), attained
by the posted unit price \(2\).
Every posted price \(p\in[1,2]\) is also globally feasible and earns
\(pS(p)=1\). Different prices give different allocations on sets of
positive measure. This establishes nonuniqueness against all feasible
mechanisms, not merely among posted prices.

\section{General finite-menu approximation}\label{sec:rounding}
The profit-discount starting point is retained from the prior catalogue
attempt~\cite{PriorAttempt}. The additional obligations here are arbitrary
nonempty compact \(Q\), signed continuous costs with a modulus, a
controlled all-type IR repair, and the effective global certificate
developed below. No priority claim is made for the general menu technique.

Fix \(0<\eta<1\), quantity precision \(\delta_q>0\), and price precision
\(\delta_p>0\).
Let \(Q_\delta\subset Q\) be a finite feasible \(\delta_q\)-net.
Let \(G\) be a scalar grid of spacing \(\delta_p\), spanning
\([-K-\eta P,K+\eta P]\) with outward-rounded endpoints.
Every value in this interval is within \(\delta_p\) of a grid price, and
\[
 |g|\le T_0:=K+\eta P+\delta_p\quad(g\in G),\qquad
 e=H\delta_q+\delta_p.
\]
For each contract \((q_j,t_j)\) of a normalized mechanism, put
\(\pi_j=t_j-C(q_j)\). Choose \(q_j'\in Q_\delta\) with
\(\norm{q_j'-q_j}\le\delta_q\), and \(g_j\in G\) within \(\delta_p\)
of \(t_j-\eta\pi_j\). Keep all distinct pairs \((q_j',g_j)\).
The result is a finite nonempty menu even when the old menu was infinite.
No baseline or outside allocation is inserted.

\begin{theorem}[All-type finite-menu loss]\label{thm:rounding}
Lower every rounded price by \(\gamma=\eta P+e+m_0\), where \(m_0\ge0\).
The resulting finite menu, with any fixed first-maximizer rule, is
globally IC and gives utility at least \(m_0\) to every type.
Its profit at each type is at least the old normalized profit minus
\(E(\delta_q,\delta_p,\eta)+m_0\), where
\begin{equation}\label{eq:rounderror}
 E(\delta_q,\delta_p,\eta)
 =2\eta P+\left(\frac2\eta-1\right)(H\delta_q+\delta_p)
             +\delta_p+\omega(\delta_q).
\end{equation}
The same expected-profit bound holds under every type distribution.
\end{theorem}
\begin{proof}
Fix a type and let \(i\) be its originally assigned contract.
For every old contract \(j\), define
\[
 U_j=\theta\cdot q_j-t_j,\qquad
 e_j=\theta\cdot(q_j'-q_j)-[g_j-(t_j-\eta\pi_j)].
\]
Then \(|e_j|\le e\), and the rounded unshifted utility is
\(U_j+\eta\pi_j+e_j\).
For a selected entry choose any old preimage \(j\) if several
contracts rounded to the same pair.
New utility maximization and old global IC yield
\[
 U_j+\eta\pi_j+e_j\ge U_i+\eta\pi_i+e_i,\qquad U_i\ge U_j.
\]
Hence \(\pi_j\ge\pi_i-2e/\eta\).
The selected unshifted profit is at least
\((1-\eta)\pi_j-\delta_p-\omega(\delta_q)\).
Its maximum utility is at least that of the rounded old assignment:
\begin{equation}\label{eq:rawIR}
 \max_j\{\theta\cdot q_j'-g_j\}
 \ge U_i+\eta\pi_i+e_i\ge-\eta P-e.
\end{equation}
The common price reduction leaves choices unchanged and gives utility
at least \(m_0\). Its pointwise profit loss is at most
\begin{align*}
 \eta P+\frac{2(1-\eta)e}{\eta}
       +\delta_p+\omega(\delta_q)+\gamma
 &=2\eta P+\left(\frac2\eta-1\right)e
       +\delta_p+\omega(\delta_q)+m_0,
\end{align*}
which is \eqref{eq:rounderror}.
Every report is assigned a maximizer from the same menu; a contract
obtained at a different report was already an available option.
This proves global IC. Equation~\eqref{eq:rawIR} proves global IR,
and every quantity remains in \(Q\).
A finite maximum of affine functions with a fixed tie rule is
measurable. No assumption about tie probabilities enters the pointwise
loss proof.
\end{proof}

For \(0<\delta<1\), \(\delta_q=\delta_p=\delta\), and
\(\eta=\sqrt\delta\), the error simplifies to
\begin{equation}\label{eq:sqrtrate}
 E=2(P+H+1)\sqrt\delta-H\delta+\omega(\delta)
 \le2(P+H+1)\sqrt\delta+\omega(\delta).
\end{equation}
Square dyadics \(\delta=2^{-2n}\) give rational \(\eta=2^{-n}\).
The approximation uses the uniform cost modulus, and requires neither
Lipschitz cost nor \(C(0)=0\). Convexity of \(C,Q\) is compatible with
this result but is not needed for this rounding argument.

\section{Effective primitives and an unconditional IR normalizer}
\label{sec:effective}
For exhaustive search, arbitrary nonempty grid subsets need not be IR.
An exact IR decision is unnecessary.
For a finite menu \(M\), let
\[
 U_M(\theta)=\max_j(\theta\cdot q_j-t_j),\qquad
 m_M=\min_\Theta U_M.
\]
Fix \(\xi>0\). Take a feasible type \(h\)-net with \(R_Qh\le\xi/2\).
At each net point compute a certified lower utility value within
\(\xi/2\) of the true value. If \(\ell\) is the smallest such value, put
\[
 \underline m=\ell-R_Qh.
\]
The Lipschitz bound and covering certificate imply
\begin{equation}\label{eq:minutility}
 m_M-\xi\le\underline m\le m_M.
\end{equation}
Indeed every type is close to a net point, and every net value is at
least \(m_M\). Lower all prices by
\(s_M=\max\{0,-\underline m\}\), and denote the new menu by \(N_\xi(M)\).
It is globally IC/IR for every \(M\), with unchanged choices.
For the menu constructed in Section~\ref{sec:rounding},
\eqref{eq:rawIR}--\eqref{eq:minutility} give
\(s_M\le\eta P+e+\xi\), so its pointwise loss is at most \(E+\xi\).
This can replace the fixed shift \(\gamma\) in the search.

Every raw grid payment is bounded by \(T_0\), while
\(U_M\ge-HR_Q-T_0\). Uniformly over all grid subsets,
\begin{equation}\label{eq:outbounds}
 s_M\le HR_Q+T_0+\xi,\qquad
 |t^{\mathrm{new}}|\le2T_0+HR_Q+\xi,\qquad
 |\pi^{\mathrm{new}}|\le P_{\mathrm{out}}:=2T_0+HR_Q+\xi+A.
\end{equation}

\subsection{An explicit effective input representation}
Compactness and continuity alone are not finite algorithmic inputs.
The computational theorem uses the following certified representation.
\begin{enumerate}
\item At every positive rational radius, finite feasible nets for
  \(Q\) and \(\Theta\), with computable-real coordinate names and valid
  covering certificates. Every point belongs exactly to its set;
  merely approximate membership in \(Q\) is insufficient.
\item The bounds \(H,R_Q,D,A\), a computable upper modulus \(\omega\),
  and certified evaluations of \(C\) at feasible named quantity points.
\item A bound \(\bar f\ge\sup_\Theta f\) and probability quadratures:
  for every \(h>0\), a finite measure
  \(\mu_h=\sum_\ell w_\ell\delta_{\theta_\ell}\), with nonnegative
  rational weights summing to one, named points in \(\Theta\), and
  certified \(W_1(\mu,\mu_h)\le h\). A coupling with mean transportation
  distance at most \(h\) suffices as the certificate.
\end{enumerate}
The input law is the original continuously positive density.
For an \(L_g\)-Lipschitz function \(g\), the coupling directly gives
\begin{equation}\label{eq:quadrature}
 \left|\int g\,d\mu-\sum_\ell w_\ell g(\theta_\ell)\right|\le L_gh.
\end{equation}
This is an ordinary continuous-function integration representation of
the probability law, not an oracle for screening, its optimum, or
discontinuous menu profit.

Price meshes and IR shifts can be rational. Quantities need not be:
an irrational singleton \(Q\) contains no rational point.
The output consists of a finite list of certified feasible coordinate
names and rational payments. The menu lets buyers choose an entry.
The fixed first-maximizer rule is a measurable mathematical direct
implementation, not a promise that exact equality at arbitrary real
reports is decidable from finite approximations.
Its potential is the explicitly specified finite maximum of affine
functions, and the global certificate requires no numerical tie decisions.
For ordinary computability the primitive oracles must themselves be
computable, with their costs included in the input model.
The analytic approximation theorem holds for every original mathematical
instance, whether or not that instance has an effective presentation.

\section{Certified integration without polynomial primitives}\label{sec:envelopes}
Fix a final finite menu with \(M\) entries and profits
\(\pi_j=t_j-C(q_j)\). Let \(P_0\ge\max_j|\pi_j|\);
\(P_{\mathrm{out}}\) in \eqref{eq:outbounds} works uniformly.
Put \(u_j(\theta)=\theta\cdot q_j-t_j\), \(V(\theta)=\max_j u_j(\theta)\).
For \(L>0\) define continuous envelopes
\begin{align*}
 g_L^+(\theta)&=\max_j\{\pi_j-L(V(\theta)-u_j(\theta))\},\\
 g_L^-(\theta)&=\min_j\{\pi_j+L(V(\theta)-u_j(\theta))\}.
\end{align*}
For every type and every selected utility maximizer \(i\),
\begin{equation}\label{eq:envelopepoint}
 g_L^-(\theta)\le\pi_i\le g_L^+(\theta).
\end{equation}
Both envelopes lie in \([-P_0,P_0]\) and are \(2LR_Q\)-Lipschitz.
For example,
\(g_L^+=\max_j(\pi_j+Lu_j)-LV\) is a difference of two
\(LR_Q\)-Lipschitz functions. The analogous minimum formula proves
the lower-envelope statement.

\begin{lemma}[Integrated envelope gap]\label{lem:envelope}
For \(b=2P_0/L\), \(\beta>0\), and
\(\Lambda_M=M(M-1)\bar f(2H)^{d-1}\),
\begin{equation}\label{eq:envelopegap}
 \int(g_L^+-g_L^-)\,d\mu
 \le H\beta+\omega(\beta)
       +b\left(1+\frac{2P_0\Lambda_M}{\beta}\right).
\end{equation}
\end{lemma}
\begin{proof}
Let \(J_b(\theta)=\{j:V(\theta)-u_j(\theta)\le b\}\).
An index outside this set has penalized upper term below \(-P_0\)
and penalized lower term above \(P_0\), whereas a true maximizer
offers an unpenalized profit in \([-P_0,P_0]\).
The envelope width is at most the profit range among \(J_b\).
For \(j,k\in J_b\), \(|u_j-u_k|\le b\) and
\[
 t_j-t_k=\theta\cdot(q_j-q_k)-(u_j-u_k).
\]
Consequently
\[
 |\pi_j-\pi_k|
 \le H\norm{q_j-q_k}+b+\omega(\norm{q_j-q_k}).
\]
If every pair in \(J_b\) is within \(\beta\) in quantity, this bounds
the width by \(H\beta+b+\omega(\beta)\).
Otherwise a pair with \(\norm{q_j-q_k}>\beta\) satisfies
\[
 |\theta\cdot(q_j-q_k)-(t_j-t_k)|\le b.
\]
This slab has thickness at most \(2b/\beta\).
Because \(\Theta\) is contained in the radius-\(H\) Euclidean ball,
each orthogonal cross-section has volume at most \((2H)^{d-1}\).
Fubini and the density bound give slab probability at most
\(2b\bar f(2H)^{d-1}/\beta\).
Union over at most \(M(M-1)/2\) pairs gives probability at most
\(\Lambda_Mb/\beta\).
On this exceptional event use width \(2P_0\).
This proves \eqref{eq:envelopegap}; no separation between distinct
quantities was assumed.
\end{proof}

For a desired profit-interval width \(\zeta>0\), choose \(\beta>0\)
with \(H\beta+\omega(\beta)\le\zeta/4\), and take rational
\[
 L\ge\max\left\{1,\frac{8P_0}{\zeta}
               \left(1+\frac{2P_0\Lambda_M}{\beta}\right)\right\}.
\]
The true integrated envelope gap is then at most \(\zeta/2\).
Use quadrature radius at most \(\zeta/(32LR_Q)\) and point evaluations
with error at most \(\zeta/16\).
Equation~\eqref{eq:quadrature} gives approximations \(I^-,I^+\) to the
two envelope integrals with absolute error at most \(\zeta/8\).
Thus
\begin{equation}\label{eq:profitinterval}
 [\,I^--\zeta/8,\ I^++\zeta/8\,]
\end{equation}
contains the exact expected menu profit and has width at most \(\zeta\).
The true gap contributes \(\zeta/2\), the two approximation errors
\(\zeta/4\), and the two outward margins \(\zeta/4\).
All numerical errors are included.

Even without a supplied density bound, increasing \(L\) and refining
integration until the certified gap is small terminates.
The envelopes converge together almost everywhere: distinct quantities
tie only on a Lebesgue-null hyperplane, while identical quantities and
payments have identical profits.
Bounded convergence proves this qualitative claim.
Lemma~\ref{lem:envelope} supplies the stronger explicit regularity
and dimension dependence.

\section{A finite exhaustive global certificate and its size}\label{sec:search}
Enumerate every nonempty subset \(M\) of \(Q_\delta\times G\).
Apply \(N_\xi(M)\) and compute a profit interval
\([\ell_M,u_M]\) of width at most \(\zeta\) by
Section~\ref{sec:envelopes}.
Every candidate is exactly globally IC/IR on all of \(\Theta\).
Let
\[
 W=\max_M J(N_\xi(M)).
\]
This finite maximum satisfies \(W\le V^*\).
Conversely, every feasible original mechanism, after the improving
normalization of Lemma~\ref{lem:normal}, maps to one of these candidates
with pointwise loss at most \(E+\xi\).
Taking its supremum gives
\begin{equation}\label{eq:searchbracket}
 W\le V^*\le W+E+\xi.
\end{equation}
This argument does not rely on attainment of the original supremum.

Choose \(\widehat M\) with the largest \(\ell_M\), using enumeration
order for ties, and define
\[
 \mathsf L=\max_M\ell_M,\qquad
 \mathsf U=\max_Mu_M+E+\xi.
\]
Since \(u_M\le\ell_M+\zeta\) for every candidate,
\begin{equation}\label{eq:globalcertificate}
 \mathsf L\le J(N_\xi(\widehat M))\le V^*\le\mathsf U,\qquad
 \mathsf U-\mathsf L\le E+\xi+\zeta.
\end{equation}
A certified rational upper bound \(\overline E\) may replace \(E\)
in these expressions.
The finite menu, its IR-net certificate, the finite objective intervals,
and \eqref{eq:globalcertificate} form the certificate.

\begin{theorem}[Certified approximation for the general compact class]
\label{thm:algorithm}
Under the effective representation of Section~\ref{sec:effective},
for every \(\varepsilon>0\) a terminating finite computation outputs a
finite menu with exact global IC/IR and a certified profit interval
whose upper endpoint exceeds its guaranteed achieved profit by at most
\(\varepsilon\).
It applies without \(0\in Q\), nonnegative cost, Lipschitz cost,
or polynomial density or cost assumptions.
\end{theorem}
\begin{proof}
Take \(\xi=\zeta=\varepsilon/4\).
Choose a sufficiently small square dyadic \(\delta\), with
\(\eta=\sqrt\delta\), so that \eqref{eq:sqrtrate} has certified upper
bound at most \(\varepsilon/2\).
The supplied modulus guarantees termination.
Apply \eqref{eq:globalcertificate}.
\end{proof}

\subsection{Cardinality and precision bounds}
A feasible quantity net can be reduced to
\begin{equation}\label{eq:netcard}
 N_Q\le(1+8R_Q/\delta_q)^d.
\end{equation}
Start with an oracle net of radius \(\delta_q/4\).
Process its points in order, approximating distances to retained
points within \(\delta_q/16\); retain a point only if every approximate
distance exceeds \(\delta_q/2\).
A discarded point is within \(9\delta_q/16\) of a retained point.
Every \(q\in Q\) is therefore within \(13\delta_q/16<\delta_q\)
of a retained point.
Retained points are more than \(7\delta_q/16\) apart.
Their disjoint balls of radius \(\delta_q/8\) lie in the ball of
radius \(R_Q+\delta_q/8\); volume comparison proves \eqref{eq:netcard}.
The argument applies also to feasible type nets and lower-dimensional \(Q\).
The price grid has at most \(2(K+\eta P)/\delta_p+3\) points.
Thus every output menu has at most
\begin{equation}\label{eq:contractcard}
 N\le(1+8R_Q/\delta_q)^d
             \left[\frac{2(K+\eta P)}{\delta_p}+3\right]
\end{equation}
entries, with at most \(2^N-1\) subsets searched.

For common mesh \(\delta\), the loss is
\(O(\sqrt\delta+\omega(\delta))\) and
\(N=O(\delta^{-(d+1)})\) at fixed primitive bounds.
If \(\omega(\delta)\le L_C\delta^\beta\), \(0<\beta\le1\), put
\(\gamma=\min\{1/2,\beta\}\).
The loss is \(O(\delta^\gamma)\); choosing
\(\delta\) of order \(\varepsilon^{1/\gamma}\) yields
\[
 N=O(\varepsilon^{-(d+1)/\gamma}),\qquad
 \text{number of menus}\le
 \exp[O(\varepsilon^{-(d+1)/\gamma})].
\]
For arbitrary continuous cost the explicit choice uses its supplied
modulus, with no universal algebraic rate asserted.
For integration, use the preceding \(L\) with \(P_0=P_{\mathrm{out}}\)
and \(M\le N\), then quadrature precision \(\zeta/(32LR_Q)\);
for IR use type radius at most \(\xi/(2R_Q)\).
These are oracle-accuracy bounds. Generating the primitive nets and
quadratures has its actual input-dependent cost; no polynomial bit
complexity follows from these counts.

\section{Concrete effective inputs with continuous primitives}\label{sec:concrete}
One concrete program-input realization takes \(\Theta,Q\) to be
nonempty rational polytopes, with \(\Theta\) full-dimensional and
\(Q\) possibly lower-dimensional or excluding zero.
Let \(C,f\) have certified evaluation programs, uniform moduli, and
bounds on their compact domains.
Convexity of \(C\), positivity of \(f\), and normalization
\(\int f=1\) are input certificates or promises, not conclusions
of numerical tests.
Neither function need be polynomial or Lipschitz.

Rational triangulations and barycentric grids give feasible rational
nets with covering radii certified by simplex diameters, also in a
proper affine hull for \(Q\). Prune them as above.
Point-coordinate operations are rational; function evaluations use
certified intervals.

The probability quadrature is constructive as well.
Cut a rational bounding cube into rational cubes of side at most
\(h/\sqrt d\), and intersect each with \(\Theta\).
The nonempty intersections are rational polytopes of diameter at most
\(h\). Select a rational feasible representative \(\theta_j\) in each,
and compute its rational volume by triangulation.
Shared boundaries are Lebesgue-null.
Evaluation of \(f\) at \(\theta_j\), with error radius \(e_f\), and
its modulus \(\omega_f(h)\) give rational cell-mass bounds
\(\ell_j,u_j\), clipping negative lower bounds to zero. They satisfy
\[
 \sum_j(u_j-\ell_j)
 \le2\vol(\Theta)\,[e_f+\omega_f(h)].
\]
They can be made arbitrarily narrow.
The true cell masses \(w_j^*\) sum to one, so rational weights
\[
 \ell_j\le w_j\le u_j,\qquad w_j\ge0,\qquad\sum_jw_j=1
\]
exist and can be obtained by successively allocating the remaining
mass within these rational intervals.
Mapping each type to its cell representative costs at most \(h\).
To change the atomic weights, retain their common masses and move
the unmatched mass, of size half the \(\ell^1\) difference, over
distance at most \(D\). Thus
\[
 W_1\left(\mu,\sum_jw_j\delta_{\theta_j}\right)
 \le h+\frac D2\sum_j|w_j-w_j^*|
 \le h+\frac D2\sum_j(u_j-\ell_j).
\]
This proves effective quadrature for general continuous \(f\) with
its modulus. Rational polynomial densities are a special case in which
cell masses can instead be integrated exactly over rational simplices.
The oracle theorem is broader than this illustration and allows any
compact convex sets with the certified representations specified above.

Quadrature support can be controlled independently of its construction.
For target precision \(\rho\), first request primitive precision
\(\rho/2\) and a feasible type net of radius \(\rho/4\).
Move each atom to an approximately nearest net point, using distance
errors at most \(\rho/16\).
The chosen distance is at most \(3\rho/8<\rho/2\).
After aggregating rational weights the total error is at most \(\rho\),
and support size is at most
\[
 (1+32H/\rho)^d.
\]
Processing the original oracle output retains its actual cost;
support compression is not a polynomial-time claim.
For example,
\[
 Q=[0,1]\times[1,2]^{d-1},\qquad C(q)=1-\sqrt{q_1}
\]
has \(0\notin Q\) and continuous convex non-Lipschitz cost, with
modulus at most \(\sqrt\delta\). It is covered by this concrete model.

\section{A sufficient finite-mass flow dual certificate}\label{sec:dual}
The following verifier applies to the full compact class.
It does not assert strong dual attainment or necessity of a finite witness.
Choose a finite positive measure \(\Lambda\) on \(\Theta\times\Theta\)
with \(\Lambda_2\ll\mu\), and require
\[
 \Lambda_1\le\mu+\Lambda_2.
\]
Supply its measurable kernel representation
\(\Lambda(dx,dy)=K_y(dx)\mu(dy)\).
Equivalently, a witness may specify such a kernel directly, with
\(\int K_y(\Theta)\,\mu(dy)<\infty\).
Define
\begin{align*}
 \nu&=\mu+\Lambda_2-\Lambda_1\ge0,\\
 v(y)&=y+\int_\Theta(y-x)K_y(dx),\\
 C_Q^*(v)&=\max_{q\in Q}\{v\cdot q-C(q)\},\\
 \mathcal D(\Lambda)&=\int_\Theta C_Q^*(v(y))\,\mu(dy).
\end{align*}
The maximum exists by compactness and continuity.
The marginal inequality also implies \(\Lambda_1\ll\mu\).
Moreover
\(\int\norm v\,d\mu\le H+2H\Lambda(\Theta\times\Theta)\) and
\(|C_Q^*(v)|\le R_Q\norm v+A\), so the displayed integrals are finite.

\begin{proposition}[Exact weak-dual residual]\label{prop:dual}
For every globally IC/IR mechanism,
\begin{align}
 \mathcal D(\Lambda)-J(U,q)
 ={}&\int[C_Q^*(v)-v\cdot q+C(q)]\,d\mu\notag\\
 &+\int[U(x)-U(y)-(x-y)\cdot q(y)]\,d\Lambda(x,y)\notag\\
 &+\int U\,d\nu.\label{eq:dualresidual}
\end{align}
All three terms are nonnegative.
\end{proposition}
\begin{proof}
Add the second and third residuals to \(J\).
The coefficient measure of \(U\) cancels:
\(-\mu+\Lambda_1-\Lambda_2+\nu=0\).
The allocation coefficient relative to \(\mu\) becomes
\[
 y-\int(x-y)K_y(dx)=v(y).
\]
The result is \(\int[v\cdot q-C(q)]\,d\mu\); subtract it from
\(\mathcal D(\Lambda)\).
Every individual feasible utility is bounded and Lipschitz on compact
\(\Theta\), up to a finite constant, so all cancellation integrals are
legitimate. Nonnegativity follows from maximization, IC, and IR.
\end{proof}
An explicit feasible witness with zero residuals certifies global
optimality. A certified residual at most \(\varepsilon\) certifies an
objective gap at most \(\varepsilon\), since every feasible profit is
at most \(\mathcal D(\Lambda)\).
Marginal inequalities, integrals, and maximization/contact conditions
must actually be checked; merely asserting existence of a witness is
not verification. Piecewise kernels with certified pushforward bounds
are one possible effective format.
No derivative of \(C\) or finite cost subgradient at every boundary
allocation is required.

\section{An exact radial optimizer against all mechanisms}\label{sec:radial}
We now give explicit structural geometry conditions in a full-dimensional
class. Radial reductions are classical; no priority claim is made
(see the historical review~\cite{Rochet2024}, Sections~3.2--3.3).
The following proof compares against every feasible mechanism, not just
radial ones.

Fix \(d\ge2\), \(R>0\), \(a>0\), \(\lambda>0\), and \(c_0\ge0\), and set
\[
 \Theta=\{\theta\in\R_+^d:\norm\theta\le R\},\qquad
 Q=\{q\in\R_+^d:\norm q\le a\},\qquad
 C(q)=c_0\norm q+\frac{\lambda}{2}\norm q^2.
\]
Types are uniform on \(\Theta\). Both sets are compact, convex, and
full-dimensional; the density is constant, continuous, and positive,
and cost is continuous and strictly convex. Here \((0,0)\) is feasible.
Write \(\Pi(q,t)=\E[t-C(q)]\).
Let \(\Omega=\{\omega\in\R_+^d:\norm\omega=1\}\) carry normalized
surface measure \(\sigma\). Polar coordinates give the type law
\(\sigma(d\omega)w(r)\,dr\), where
\[
 w(r)=\frac{dr^{d-1}}{R^d},\qquad W(r)=\left(\frac rR\right)^d.
\]
Its radial virtual value is
\begin{equation}\label{eq:virtual}
 \phi(r)=r-\frac{1-W(r)}{w(r)}
 =\frac{d+1}{d}r-\frac{R^d}{dr^{d-1}},\qquad 0<r\le R.
\end{equation}
It is strictly increasing, with
\[
 \phi'(r)=\frac{d+1}{d}+\frac{(d-1)R^d}{dr^d}>0,\qquad
 \lim_{r\downarrow0}\phi(r)=-\infty,\quad \phi(R)=R.
\]
Define
\begin{equation}\label{eq:radialcandidate}
 \begin{aligned}
 \rho(r)&=\min\left\{a,\max\left\{0,\frac{\phi(r)-c_0}{\lambda}\right\}\right\},
          &&r>0,\qquad \rho(0)=0,\\
 u(r)&=\int_0^r\rho(s)\,ds,& U_*(\theta)&=u(\norm\theta),\\
 q_*(\theta)&=\rho(\norm\theta)\frac{\theta}{\norm\theta}
          &&(\theta\ne0),\qquad q_*(0)=0,\\
 t_*(\theta)&=\norm\theta\,\rho(\norm\theta)-u(\norm\theta).
 \end{aligned}
\end{equation}
The function \(\rho\) is continuous, nonnegative, nondecreasing, and
vanishes on a neighborhood of zero.

\begin{theorem}[Radial optimum]\label{thm:radial}
The mechanism \eqref{eq:radialcandidate} is globally IC/IR and maximizes
profit over all measurable globally feasible mechanisms.
Its value is
\begin{equation}\label{eq:radialvalue}
 V_*=\int_0^Rw(r)\left[(\phi(r)-c_0)\rho(r)
                         -\frac{\lambda}{2}\rho(r)^2\right]\,dr.
\end{equation}
Its utility is the unique optimal utility everywhere. Allocation and
payment are unique almost everywhere, and in fact at every interior type.
Boundary allocation uniqueness is not asserted.
\end{theorem}

\subsection{An upper bound for every globally feasible competitor}
For any such competitor, the two IC inequalities give an
\(a\)-Lipschitz convex utility with \(U(0)\ge0\).
For each \(\omega\), \(v_\omega(r)=U(r\omega)\) is convex and Lipschitz,
hence absolutely continuous, with
\begin{equation}\label{eq:rayutility}
 v_\omega(r)=U(0)+\int_0^r b_\omega(s)\,ds,\qquad
 b_\omega(r)=q(r\omega)\cdot\omega\quad\text{a.e. in }r.
\end{equation}
Indeed global IC restricted to the ray makes \(q(r\omega)\cdot\omega\)
a one-dimensional subgradient, hence the derivative at each
differentiability point. This uses no ambient differentiability theorem.
The exceptional set is one-dimensional null for each ray, and Fubini
applies because \(q\) is bounded and measurable.
Nonnegativity and Cauchy--Schwarz give
\[
 0\le b_\omega(r)\le\norm{q(r\omega)}\le a.
\]
All integrals below are finite: \(U,q\) are bounded for an individual
mechanism, and \(w(r)|\phi(r)|\) is integrable at zero.
Using
\[
 \int_0^Rw(r)\int_0^r b(s)\,ds\,dr
 =\int_0^R[1-W(s)]b(s)\,ds
\]
in \eqref{eq:rayutility} gives the exact identity
\begin{equation}\label{eq:radialprofitidentity}
 \Pi(q,t)=-U(0)+
 \int_\Omega\int_0^Rw(r)
 \left[\phi(r)b_\omega(r)-c_0\norm{q(r\omega)}
                  -\frac{\lambda}{2}\norm{q(r\omega)}^2\right]
 dr\,\sigma(d\omega).
\end{equation}
The only utility anchor is \(U(0)\); there is no omitted boundary term.

For \(b\in[0,a]\), let
\[
 g_r(b)=(\phi(r)-c_0)b-\frac{\lambda}{2}b^2.
\]
Its unique maximizer is \(\rho(r)\).
Since \(c_0\ge0\) and \(\norm q\ge b\ge0\),
\begin{equation}\label{eq:radialupper}
 \Pi(q,t)
 \le-U(0)+\int_\Omega\int_0^Rw(r)g_r(b_\omega(r))\,dr\,d\sigma
 \le V_*.
\end{equation}
This bound holds for every globally IC/IR mechanism.

\subsection{Attainment and an exact deficit identity}
The scalar \(u\) is convex and nondecreasing.
For \(\theta,\theta'\) and \(s\in[0,1]\), the triangle inequality and
scalar convexity give
\[
 u(\norm{s\theta+(1-s)\theta'})
 \le u(s\norm\theta+(1-s)\norm{\theta'})
 \le su(\norm\theta)+(1-s)u(\norm{\theta'}).
\]
Extend \(\rho\) continuously and nondecreasingly beyond \(R\), for
example constantly. Then \(U_*\) is convex on all of \(\R^d\) and
continuously differentiable, including at the origin since \(\rho\)
vanishes in a neighborhood there.
Its gradient on \(\Theta\) is \(q_*\in Q\).
The global supporting-gradient inequality proves IC at every type,
including thresholds and boundary points, for
\(t_*=\theta\cdot q_*-U_*\).
Also \(U_*\ge0\) and \(U_*(0)=0\), proving IR.
Payments are nonnegative because
\(\int_0^r\rho(s)\,ds\le r\rho(r)\), although transfer nonnegativity
was not assumed. For this mechanism \(b=\norm{q_*}=\rho\), so both
bounds in \eqref{eq:radialupper} are equalities.
This proves optimality and \eqref{eq:radialvalue}.

There is a stronger certificate. For \(r>0\), put
\[
 b=q(r\omega)\cdot\omega,\qquad
 e(r)=\phi(r)-c_0-\lambda\rho(r).
\]
Quadratic algebra gives
\[
 g_r(\rho)-g_r(b)
 =\frac{\lambda}{2}(\rho-b)^2+e(r)(\rho-b).
\]
If \(\rho=0\), then \(e\le0\) and \(\rho-b\le0\);
if \(0<\rho<a\), then \(e=0\);
if \(\rho=a\), then \(e\ge0\) and \(\rho-b\ge0\).
Thus the last product is always nonnegative.
The radial decomposition yields
\[
 \norm{q-\rho\omega}^2=(b-\rho)^2+\norm q^2-b^2.
\]
Subtracting \eqref{eq:radialprofitidentity} from \(V_*\) now gives
\begin{equation}\label{eq:deficit}
 V_*-\Pi(q,t)=U(0)+
 \E_{r,\omega}\left[
   \frac{\lambda}{2}\norm{q(r\omega)-\rho(r)\omega}^2
   +c_0(\norm{q(r\omega)}-b)+e(r)(\rho(r)-b)\right].
\end{equation}
Every term is nonnegative. The virtual coefficient is defined only
for \(r>0\); its value at the null point zero is immaterial.
Integrability follows from the same \(w|\phi|\) bound as before.

Equality of profit forces \(U(0)=0\) and \(q=q_*\) almost everywhere.
For almost every direction, Fubini and \eqref{eq:rayutility} imply
\(U(r\omega)=u(r)\) for every \(r\).
The resulting full-measure set is dense in \(\Theta\), and both
utilities are continuous, so \(U=U_*\) everywhere.
Payments agree almost everywhere. Conversely these conclusions for
a globally feasible mechanism give equal profit.
This proves the uniqueness assertions of Theorem~\ref{thm:radial}.
Moreover \(U_*\) is differentiable at every interior type, so its
relative subdifferential there is the singleton \(\{\nabla U_*\}\).
Global IC therefore forces every optimal allocation to equal \(q_*\)
at every interior type, not just almost everywhere.
This pointwise fact justifies transfer of the classical rank statements
below to all optimizers; almost-everywhere equality alone would not.

For every feasible mechanism with profit at least \(V_*-\varepsilon\),
\eqref{eq:deficit} also gives
\begin{equation}\label{eq:stability}
 U(0)\le\varepsilon,\qquad
 \E\norm{q-q_*}^2\le\frac{2\varepsilon}{\lambda}.
\end{equation}
Neither certificate is inferred from a numerical type grid.

\section{Exact exclusion and served rank-deficient bunching}\label{sec:geometry}
If \(c_0\ge R\), then \(\phi(r)\le R\le c_0\), so the displayed
optimizer has \(\rho=0\) throughout.
Suppose \(c_0<R\). There is a unique \(r_0\in(0,R)\) with
\(\phi(r_0)=c_0\). The zero-allocation region of \(q_*\) is exactly
\(\{\theta:\norm\theta\le r_0\}\), of probability
\((r_0/R)^d\in(0,1)\); exclusion is partial.
When \(c_0=0\),
\[
 r_0=\frac R{(d+1)^{1/d}},\qquad
 \Prob\{\text{exclusion}\}=\frac1{d+1}.
\]
For \(c_0>0\) the radius is strictly larger, increasing with \(c_0\).

If \(c_0+\lambda a<R\), a unique \(r_1\in(r_0,R)\) solves
\(\phi(r_1)=c_0+\lambda a\), and
\[
 \rho(r)=
 \begin{cases}
 0,&0\le r\le r_0,\\
 (\phi(r)-c_0)/\lambda,&r_0<r<r_1,\\
 a,&r_1\le r\le R.
 \end{cases}
\]
Every threshold satisfying \(\phi(r)=z\) is equivalently the unique
positive root of
\begin{equation}\label{eq:thresholdpoly}
 (d+1)r^d-dzr^{d-1}-R^d=0.
\end{equation}
Strict increase of \(\phi\) proves uniqueness; for \(z<R\) its
endpoint signs put the root in \((0,R)\).
These are primitive conditions and scalar algebraic certificates.

For \(r>0\) away from thresholds, write \(\omega=\theta/r\).
Differentiation gives
\begin{equation}\label{eq:rankmatrix}
 Dq_*(\theta)=\rho'(r)\omega\omega^\T
       +\frac{\rho(r)}r(I-\omega\omega^\T).
\end{equation}
The radial eigenvalue is \(\rho'(r)\), and the \(d-1\) tangential
eigenvalues are \(\rho(r)/r\).
On \(r_0<r<r_1\), both are positive, so the rank is \(d\).
On the outer open annulus \(r_1<r<R\), the radial eigenvalue is zero
and the tangential eigenvalues equal \(a/r>0\).
Its rank is \(d-1\) and its probability is \(1-(r_1/R)^d>0\).
Angular-interior types receive strictly positive quantities in all
coordinates. Along a fixed direction every radius in the annulus
receives \(a\omega\), but different directions receive different
allocations. A particular positive allocation has a preimage in one
ray, of \(d\)-dimensional measure zero because \(d\ge2\).
This is served rank-deficient bunching, not whole-region pooling or
a finite-menu artifact. The excluded inner region separately has rank zero.

If \(c_0<R\) but \(c_0+\lambda a\ge R\), no open saturation annulus
exists. Every served interior type has full rank; equality saturates
only at the null outer boundary.
Thus, within this class, a positive-measure served rank-\((d-1)\)
annulus exists if and only if
\begin{equation}\label{eq:bunchcondition}
 0\le c_0<R,\qquad 0<a<\frac{R-c_0}{\lambda}.
\end{equation}
For \(R=\lambda=1\), \(c_0=0\), and \(0<a<1\), partial exclusion,
a separating region, and served bunching coexist for every \(d\ge2\).
The phenomenon occurs where the declared radial quantity cap binds;
no claim about unconstrained bunching is inferred.

\section{Certified rational computation for the radial family}\label{sec:radialcompute}
This is a finite-description computation for the entire radial family,
independent of the general exhaustive menu search.
Allow also the degenerate case \(a=0\).
If \(c_0\ge R\) or \(a=0\), the zero mechanism is exactly optimal.
Otherwise put
\[
 a_*=\min\{a,(R-c_0)/\lambda\}>0,\qquad
 \delta=\frac{a_*}{N},
\]
where \(N\) is a positive integer.
For \(j=1,\ldots,N\), let \(r_j\in(0,R]\) be the unique root of
\begin{equation}\label{eq:hingeroot}
 \phi(r_j)=c_0+\lambda j\delta,\qquad
 (d+1)r_j^d-d(c_0+\lambda j\delta)r_j^{d-1}-R^d=0.
\end{equation}
These roots are strictly increasing. The last is \(R\) when the
quantity cap does not bind on a positive-measure region.
Define
\begin{align*}
 h_N(r)&=\delta\sum_{j=1}^N(r-r_j)_+,&
 \rho_N(r)&=\delta\,\#\{j:r_j\le r\},\\
 q_N(\theta)&=\rho_N(\norm\theta)\frac{\theta}{\norm\theta}
       \quad(\theta\ne0),&
 q_N(0)&=0,\\
 t_N(\theta)&=\theta\cdot q_N(\theta)-h_N(\norm\theta).
\end{align*}
The hinge function is convex and nondecreasing, with slopes in
\([0,a_*]\). Its composition with the norm is convex.
At a threshold the selected slope is the right derivative, hence
a valid scalar subgradient. If \(r>0\), a nonnegative scalar support
slope \(b\) gives the vector support slope \(b\theta/r\), because
\(\norm{\theta'}\ge(\theta/r)\cdot\theta'\).
At zero the hinge vanishes nearby.
Thus the supporting-plane inequality proves global IC at every type,
including all thresholds and the outer boundary.
The utility is nonnegative and the quantity lies in \(Q\).

Away from thresholds, \(\rho_N\) rounds \(\rho\) down in steps \(\delta\).
Where \(0<\rho<a\), the coefficient \(e(r)\) in \eqref{eq:deficit}
vanishes and \(|\rho-\rho_N|<\delta\).
In the excluded region both are zero.
In the original cap-saturated region both equal \(a\), because
\(a_*=a\) and the last threshold has been crossed.
Both allocations are radial, so the \(c_0\) term in \eqref{eq:deficit}
vanishes. Hence
\begin{equation}\label{eq:hingeerror}
 0\le V_*-\Pi(q_N,t_N)
 =\frac{\lambda}{2}\int_0^Rw(r)|\rho(r)-\rho_N(r)|^2\,dr
 \le\frac{\lambda a_*^2}{2N^2}.
\end{equation}
This bound is relative to every feasible continuum mechanism.

For rational output, assume \(R,c_0,\lambda,a\) are rational.
Fix \(m\ge1\) and \(h=R/2^m\).
By binary search on \(\{kh:0\le k\le2^m\}\), find the first grid
point \(r_j^+\ge r_j\), using the rational polynomial in
\eqref{eq:hingeroot}.
Its sign agrees with \(\phi(r)-(c_0+\lambda j\delta)\) for \(r>0\);
strict increase of \(\phi\) therefore identifies the unique crossing,
including equality. We have
\[
 0\le r_j^+-r_j<h.
\]
The rounded thresholds are nondecreasing, and repeated ones are allowed.
Replacing each \(r_j\) by \(r_j^+\) leaves the same convex-composition
proof of exact IC/IR and quantity feasibility valid.

Write \(b_N,b_N^+\) for the two radial allocations. Then
\[
 b_N-b_N^+=\delta\sum_j\ind_{[r_j,r_j^+)}\ge0.
\]
As \(w(r)\le d/R\),
\[
 \int_0^Rw(r)|b_N-b_N^+|\,dr\le\frac{a_*dh}{R}.
\]
Where they differ, \(r\) is at least the first hinge threshold \(r_1\),
\(\phi(r)\ge c_0\), and
\(\phi(r)\le R\).
For \(b,b'\in[0,a_*]\), the virtual-surplus difference is
\[
 \left[(\phi-c_0)b-\frac{\lambda b^2}{2}\right]
 -\left[(\phi-c_0)b'-\frac{\lambda(b')^2}{2}\right]
 =(b-b')\left[\phi-c_0-\frac{\lambda(b+b')}{2}\right].
\]
Its absolute value on those intervals is at most
\((R+c_0+\lambda a_*)|b-b'|\).
Combining with \eqref{eq:hingeerror} gives the certified global gap
\begin{equation}\label{eq:rationalgap}
 0\le V_*-\Pi(q_N^+,t_N^+)
 \le B_{N,m}:=\frac{\lambda a_*^2}{2N^2}
       +(R+c_0+\lambda a_*)\frac{a_*dh}{R}.
\end{equation}
Choose \(N,m\) to make the two terms at most \(\varepsilon/2\)
each. This terminates; root isolation uses at most \(N(m+1)\)
rational polynomial sign evaluations and the output has \(N\) hinges.

The achieved profit is exactly computable by rational arithmetic.
Between consecutive rounded thresholds the amount \(b\) is constant.
Since \(w(r)\phi(r)=(d+1)r^d/R^d-1\), its profit integrand is
\[
 b\left[\frac{(d+1)r^d}{R^d}-1\right]
 -\left(c_0b+\frac{\lambda b^2}{2}\right)\frac{dr^{d-1}}{R^d},
\]
with antiderivative
\[
 b\left[\frac{r^{d+1}}{R^d}-r\right]
 -\left(c_0b+\frac{\lambda b^2}{2}\right)\frac{r^d}{R^d}.
\]
Sum over at most \(N+1\) intervals, assigning repeated endpoints
zero-length intervals.
The rational achieved value \(J_{N,m}\) and \eqref{eq:rationalgap}
give the interval \([J_{N,m},J_{N,m}+B_{N,m}]\) containing the
unrestricted optimum. The curved type domain and its volume constant
do not obstruct this exact computation.

For example, \(d=2\), \(R=\lambda=1\), \(c_0=0\), \(a=1/2\), and
\(\varepsilon=1/100\) admit \(N=5\), \(\delta=1/10\), \(m=9\),
and \(h=1/512\), with thresholds
\[
 \frac{157}{256},\quad\frac{83}{128},\quad\frac{11}{16},
 \quad\frac{93}{128},\quad\frac{197}{256}.
\]
Rational evaluation gives
\[
 J_{5,9}=\frac{9890879}{83886080},\qquad
 B_{5,9}=\frac{203}{25600}<\frac1{100},\qquad
 V_*\in\left[\frac{9890879}{83886080},
             \frac{52780347}{419430400}\right].
\]
The accompanying \texttt{radial\_certificate.py} reproduces this
illustration of the proved bound; it does not replace the proof.
Approximate hinges may have additional flat radial pieces.
Their geometry is not the geometry of the exact optimizer, which was
proved independently in Section~\ref{sec:geometry}.

\section{Scope, provenance, and the certification agenda}\label{sec:scope}
The general result concerns the original compact convex quantity class,
including lower-dimensional \(Q\), \(0\notin Q\), and signed continuous
convex costs without a Lipschitz bound.
Theorem~\ref{thm:existence} gives an actual optimizer;
Theorem~\ref{thm:unique} states the precise additional condition for
utility uniqueness and almost-everywhere allocation uniqueness.
Theorem~\ref{thm:algorithm} supplies globally feasible finite menus
and certified objective gaps under explicit effective inputs, with
dimension, modulus, and oracle-accuracy dependence.
An unspecified continuous primitive is not silently treated as
a computable finite input.

Proposition~\ref{prop:dual} is a sufficient global verifier.
It does not assert that every optimizer has an attained finite-mass
flow witness or that a separate strong-duality theorem has been proved.
A certified optimizer with a verified zero region or rank-deficient
region establishes genuine optimizer geometry.
To transfer a classical derivative-rank statement to all optimizers,
almost-everywhere equality alone is insufficient; differentiability
of the common optimal utility on the region supplies pointwise
uniqueness there. Section~\ref{sec:radial} verifies this stronger fact
for the exact radial solution.

Theorems~\ref{thm:radial} and \eqref{eq:bunchcondition} give explicit
structural conditions for partial exclusion, separation, and served
rank-\((d-1)\) bunching in a declared full-dimensional class.
They do not purport to give necessary and sufficient geometry conditions
for every arbitrary density and convex cost.
The exact deficit identity and rational hinge computation give
additional all-competitor certificates for that class.
Finite-menu or hinge approximation geometry is not substituted for
the geometry of the exact optimizer.

The prior catalogue attempt~\cite{PriorAttempt}, at official revision
\texttt{00ced55}, already proved a cost-sensitive profit-discount
rounding argument for a Lipschitz-cost subclass containing the outside
allocation, with polynomial/polytope integration.
The present manuscript explicitly retains and attributes that starting
argument. It supplies the additional general-\(Q\) IR repair and
continuous-cost modulus, effective continuous-envelope integration and
global search certificate, all-boundary compact existence proof,
precise uniqueness and nonuniqueness results, a sufficient flow
verifier, and the exact radial geometry and computation module.
No novelty or priority claim is made for profit discounting, menu
approximation, radial reduction, or the broader screening literature.

Rochet's review~\cite{Rochet2024} gives the historical and structural
context, including Sections~3.2--3.3.
Carlier, Dupuis, Rochet, and Thanassoulis~\cite{CDRT2024} provide
finite-type algorithmic context, while
Babaioff, Gonczarowski, and Nisan~\cite{BGN} provide menu-size context.
These references are not represented as proofs of this manuscript's
exact collection of assumptions and certificates.
Every mathematical guarantee used here has been given its argument.

The manuscript consists of hand proofs. Numerical examples illustrate
proved inequalities; they do not replace them.
It makes no Lean formalization, external-peer-review, or novelty claim.
The complete hand proof and its transcription passed internal mathematical
and scope review. The solution claim remains subject to community review.

\begin{thebibliography}{9}
\bibitem{Rochet2024}
Jean-Charles Rochet.
\newblock \emph{Multidimensional Screening after 37 Years}.
\newblock 2024. TSE Working Paper 24-1536; published in
\emph{Journal of Mathematical Economics} \textbf{113}, article 103010.
\newblock \href{https://www.tse-fr.eu/sites/default/files/TSE/documents/doc/wp/2024/wp_tse_1536.pdf}
{Primary review PDF}. Relevant historical discussion: Sections 3.2--3.3.

\bibitem{CDRT2024}
Guillaume Carlier, Xavier Dupuis, Jean-Charles Rochet, and John Thanassoulis.
\newblock A general solution to the quasi linear screening problem.
\newblock \emph{Journal of Mathematical Economics} \textbf{114} (2024),
article 103025.
\newblock \href{https://xdupuis.perso.math.cnrs.fr/publications/screening-algo_paper.pdf}
{Published article PDF}.
\newblock \href{https://doi.org/10.1016/j.jmateco.2024.103025}
{doi:10.1016/j.jmateco.2024.103025}.

\bibitem{BGN}
Moshe Babaioff, Yannai A. Gonczarowski, and Noam Nisan.
\newblock \emph{The Menu-Size Complexity of Revenue Approximation}.
\newblock \href{https://arxiv.org/abs/1604.06580v4}{arXiv:1604.06580v4}.

\bibitem{PriorAttempt}
Prior catalogue research attempt, attributed to GPT 6 Astra Ultra.
\newblock \emph{Multidimensional screening: a globally
incentive-compatible finite-menu certificate}.
\newblock \texttt{gpt\_6\_astra\_ultra/D82-92.tex},
official repository revision \texttt{00ced55}, first posted
9 October 2026; marked unresolved.
\newblock \href{https://github.com/mehmetmars7/Erdosproblems-llm-hunter/blob/00ced55ba1223fb48603e1f15028d753997ee515/attacks/open_problems/economics/gpt_6_astra_ultra/D82-92.tex}
{Pinned source at the cited revision}.
\newblock Source of the retained cost-sensitive profit-discount
rounding argument.
\end{thebibliography}
\end{document}
